\documentclass[11pt]{article}

\usepackage[margin=1in]{geometry}
\usepackage{amsmath}
\usepackage{amssymb}
\usepackage{booktabs}
\usepackage[hidelinks]{hyperref}

\title{Hat layout in WoolyGraphs:\\crown shaping and edge-length uniformity}
\author{}
\date{2 September 2026}

\begin{document}
\maketitle

\begin{abstract}
Notes on two questions about the circular-hat pipeline. First, what makes a
knitted crown close like a sphere, and how well do our charts do it. Second,
how uniform are the edge lengths the layout optimizer produces, measured
against the gauge they are supposed to hit. All measurements are at
8 stitches per inch horizontal and 12 rounds per inch vertical
gauge. These are single-run measurements of the current code, recorded as a
baseline for comparing optimization protocols later, not averaged trials.
\end{abstract}

\section{Setup and notation}

A chart is one pattern repeat: rounds in knit order, one cell per stitch slot,
blank where no stitch exists. The hat is that chart tiled $P$ times around.
Write $h$ for the horizontal gauge in stitches per inch, $v$ for the vertical
gauge in rounds per inch, $W_0$ for the live stitches per repeat at the widest
round, and $N_0 = P W_0$ for the total.

The layout is a graph on stitch positions in $\mathbb{R}^3$. Two edge families
carry the gauge:
\begin{itemize}
  \item \textbf{horizontal} (ring): each stitch to its right-hand neighbour in
        the same round, rest length $1/h = 0.125\,\text{in}$;
  \item \textbf{vertical} (column): each stitch to the stitch above it in its
        column, rest length $1/v \approx 0.0833\,\text{in}$.
\end{itemize}
Column edges with a decrease or a cast-on at either end are counted separately
as \textbf{shaping} edges. They nominally target $1/v$ as well, but a
\texttt{k2tog} merges two columns into one, so there is no reason for them to
sit at plain column gauge, and mixing them into the vertical bucket smears it.

The initial layout places round $r$, stitch $i$ of $n_r$ on a helix at
fractional turn $r + i/n_r$, at the radius whose chord between neighbouring
stitches is exactly $1/h$. The cast-on round is flat at $z=0$ and is held fixed
during optimization, so the fabric hangs from a brim that is gauge-exact by
construction.

\section{Crown shaping}

\subsection{The law}

A hat closes like a sphere when each round's live stitch count is the
circumference of the sphere at that latitude. Taking the sphere radius from the
brim,
\begin{equation}
  R = \frac{N_0}{2\pi h},
  \qquad
  W(r) = W_0 \cos\!\left(\frac{r - r_{\text{eq}}}{v R}\right),
  \label{eq:sphere}
\end{equation}
where $r_{\text{eq}}$ is the round at which the cylinder ends. Rounds advance
by arc length along the surface, $1/v$ per round, so the angle from the equator
after $r - r_{\text{eq}}$ rounds is $(r-r_{\text{eq}})/(vR)$. The crown
therefore needs
\begin{equation}
  \frac{\pi}{2} v R
  \label{eq:height}
\end{equation}
rounds to reach the pole. This is the same relation as the
\texttt{Sphere Hat Decs} spreadsheet, and it is what the Shaping tab overlays
on the chart.

Equation~\eqref{eq:height} is the constraint that bites in practice. It couples
three things a chart author sets independently: the repeat count, the gauge,
and how many rounds of chart sit above the cylinder. A crown that is too short
for its own circumference cannot be fixed by moving decreases around.

\subsection{Method}

To ask whether a chart's crown \emph{is} a sphere, fit
equation~\eqref{eq:sphere} to the chart's live-stitch profile with
$r_{\text{eq}}$ as the only free parameter, and report the root-mean-square
residual in stitches per repeat.

The equator must be fitted rather than read off the chart. The obvious
estimator, ``the last round at full width'', is biased: near the equator the
cosine is flat, so a correctly shaped chart holds full width for several rounds
past the true $r_{\text{eq}}$, and taking the last of them compresses the
remaining crown and manufactures a large residual. On the reshaped
\texttt{alt\_cubes} below, that estimator puts the equator seven rounds too
high and reports a residual ten times the fitted one.

\subsection{Results}

\begin{table}[h]
\centering
\begin{tabular}{lrrrrrr}
\toprule
chart & $P$ & cast-on & circ.\ (in) & crown rounds & needed & rms \\
\midrule
\texttt{small\_cubes}            & 4 & 136 & 17.0 & 50 & 54.0 & 1.41 \\
\texttt{alt\_cubes}, as extracted & 6 & 120 & 15.0 & 50 & 49.5 & 0.74 \\
\texttt{alt\_cubes}, reshaped    & 6 & 120 & 15.0 & 49 & 49.5 & 0.29 \\
\texttt{alt\_cubes}, reshaped    & 7 & 140 & 17.5 & 55 & 57.7 & 0.76 \\
\bottomrule
\end{tabular}
\caption{Best-fit sphere residual, in stitches per repeat. ``Crown rounds'' is
what the chart provides above the fitted equator; ``needed'' is
equation~\eqref{eq:height}.}
\label{tab:shape}
\end{table}

\texttt{small\_cubes} is the reference: at four repeats it tracks the sphere to
$\pm 2.5$ stitches per repeat, or about 4\% of its repeat width, and it is the
chart whose knitted result is known to be good. Its residual is not the
smallest in the table, which is worth keeping in mind before treating small
residuals as the goal in themselves.

The extracted \texttt{alt\_cubes} crown fits at rms 0.74, but the residual
understates two visible defects. It holds full width about ten rounds too long
and then decreases faster than the sphere to catch up, and it reaches one
stitch per repeat five rounds before the chart ends, so the last five rounds
are a single-stitch spike rather than a closure. Reshaping re-times the same
gore geometry onto equation~\eqref{eq:sphere}, keeping the colours and the
column at which the wedge opens: the profile then agrees with the sphere to
within $\pm 0.7$ stitches per repeat and closes smoothly ($5,5,4,3,2,2,1,1$
over the last eight rounds, against $5,3,3,1,1,1,1,1$ before).

The shaping in a chart is correct for one repeat count only. At $P=7$ the same
reshaped chart is a 17.5\,\text{in} hat, closer to \texttt{small\_cubes} in size,
but equation~\eqref{eq:height} then asks for 58 crown rounds where the chart has
55, and the residual roughly triples. That is still better than the chart
started, so $P=7$ is a defensible trade of shape for size; making it exact means
re-running the reshape at $P=7$ and adding about three rounds of crown.

\section{Edge-length uniformity}

\subsection{Method}

The layout optimizer minimises a weighted sum of three terms: a gauge term
(squared relative error of every edge against its rest length), an inflate term
($1/r$ repulsion from the centroid, which keeps the surface from collapsing),
and a smoothness term on the solid angle subtended at each vertex by its
neighbour fan.

To measure the result, take the same edge set the optimizer works on and record
the realised length of every edge. For \texttt{small\_cubes} at four repeats
that is 16{,}888 edges: 8{,}448 horizontal, 8{,}016 vertical, and
424 shaping. Report per-class mean, median and standard deviation of the
lengths, and summarise the whole layout by the root-mean-square of the relative
error $\ell_e / \ell^{\,*}_e - 1$.

\subsection{Results}

\begin{table}[h]
\centering
\begin{tabular}{llrrrr}
\toprule
layout & class & gauge (in) & mean (in) & std (in) & mean vs gauge \\
\midrule
initial helix
 & horizontal & 0.1250 & 0.1252 & 0.0023 & $+0.2\%$ \\
 & vertical   & 0.0833 & 0.0996 & 0.0365 & $+19.5\%$ \\
 & shaping    & 0.0833 & 0.1517 & 0.0504 & $+82.0\%$ \\
\addlinespace
after 1000 iterations
 & horizontal & 0.1250 & 0.1279 & 0.0063 & $+2.3\%$ \\
 & vertical   & 0.0833 & 0.0904 & 0.0056 & $+8.5\%$ \\
 & shaping    & 0.0833 & 0.0990 & 0.0084 & $+19.9\%$ \\
\bottomrule
\end{tabular}
\caption{Edge lengths for \texttt{small\_cubes} at four repeats, before and
after 1000 Adam iterations at learning rate 0.002 with weights
$\lambda_{\text{gauge}}=1.0$, $\lambda_{\text{inflate}}=0.02$,
$\lambda_{\text{smooth}}=0.1$.}
\label{tab:edges}
\end{table}

Overall relative error falls from 36.8\% to 9.1\% rms. Three observations.

\paragraph{The helix is exact where it is trivially exact.}
Horizontal edges start at $+0.2\%$ with a standard deviation of
0.0023\,\text{in}, because the helix places each ring at the radius that makes
the chord exactly $1/h$. Vertical edges are exact in the cylinder for the same
reason and wrong only in the crown, which is why their initial median is
exactly $1/v$ while their mean is $+19.5\%$. A mean alone would have suggested
the whole fabric was stretched; the median says the error is confined to the
shaping region. Shaping edges start $+82\%$ long, since a decrease edge spans
two columns in the initial placement.

\paragraph{Optimization trades peak error for spread.}
Root-mean-square error improves by a factor of four, and the vertical standard
deviation tightens by a factor of six, from 0.0365 to
0.0056\,\text{in}. But the number of edges worse than 10\% off gauge
\emph{rises}, from 1{,}761 to 3{,}010: the optimizer removes a small
population of very large errors and distributes a modest error widely. Any
comparison between protocols should therefore report a distribution statistic,
not a threshold count, which can move the wrong way while the layout improves.

\paragraph{There is a systematic stretch.}
Every class ends up long: $+2.3\%$ on rings, $+8.5\%$ on columns, $+19.9\%$ on
shaping edges. The signs agree across classes and the effect survives further
iterations, so this is the inflate term pulling against the gauge term at
equilibrium rather than incomplete convergence. Lowering
$\lambda_{\text{inflate}}$ should move all three distributions back onto their
targets, and how far it can be lowered before the surface collapses is the
obvious next experiment.

\subsection{Tooling}

The Edges tab plots the three distributions overlaid on shared bins with the
gauge for each drawn as a rule, and lists mean, median and standard deviation
beneath. The count axis is logarithmic by default: these distributions put
thousands of edges in one bin and tens in the tail, and the tail is the part
the question is about. The x axis can switch from inches to error relative to
gauge, which places all three classes on a single reference line at zero.

The strain heatmap in the main view replaces the stitches with their edges,
coloured by signed error against gauge: black on gauge, orange too long, teal
too short. On-gauge edges recede into the background, so what remains visible
is exactly the problem regions. Both views read the same edge set as the
optimizer, so they cannot disagree with it.

\section{Caveats and next steps}

Every number here is one run of one configuration, on charts of roughly
8{,}500 stitches. Nothing is averaged over seeds, and the optimizer's
starting point is deterministic, so run-to-run variation is unmeasured.

The rest length assigned to shaping edges is a modelling choice, not a
measurement. Calling it $1/v$ makes every decrease edge look 20\% long even in
a good layout. A \texttt{k2tog} spans two columns at the round below, so a rest
length nearer $\sqrt{(1/v)^2 + (1/2h)^2} \approx 0.104\,\text{in}$ may be the
honest target; adopting it would change the shaping row of
Table~\ref{tab:edges} substantially and is worth settling before protocols are
compared on that class.

For comparing optimization protocols, the measurements to hold fixed are the
chart, the repeat count, the gauge, and the iteration budget; the statistics to
report are rms relative error overall and per-class standard deviation. The
current method is the baseline in Table~\ref{tab:edges}.

\end{document}
